\documentclass{article}

% ============================================================
% MA 213 In-Class Worksheet Template
% Student-facing worksheet for lecture activities.
% ============================================================

\usepackage[margin=1in]{geometry}
\usepackage{xcolor}
\usepackage{parskip}
\usepackage{array}
\usepackage{enumitem}
\usepackage{listings}

\definecolor{codegray}{gray}{0.97}
\definecolor{huddleblue}{RGB}{214,234,255}
\definecolor{predictionyellow}{RGB}{255,242,200}
\definecolor{debatepink}{RGB}{255,225,225}

\lstset{
  basicstyle=\ttfamily\small,
  columns=fullflexible,
  frame=single,
  framerule=0.4pt,
  backgroundcolor=\color{codegray},
  breaklines=true,
  showstringspaces=false,
  keepspaces=true,
  xleftmargin=0pt,
  xrightmargin=0pt
}

\newcommand{\coursenumber}{MA 213}
\newcommand{\weekplaceholder}{2}
\newcommand{\lectureplaceholder}{3}
\newcommand{\lecturetitle}{Sampling and Bias: Think/Pair/Share}

\newcommand{\nameblank}{\rule{1.75in}{0.4pt}}
\newcommand{\idblank}{\rule{1.55in}{0.4pt}}
\newcommand{\shortblank}{\rule{1.6in}{0.4pt}}
\newcommand{\longblank}{\rule{0.93\linewidth}{0.4pt}}

\newcommand{\responsebox}[2]{
  \noindent\fbox{%
    \begin{minipage}{0.95\linewidth}
      \textbf{#1} #2
      \vspace{0.55in}
    \end{minipage}
  }
}

\newcommand{\linedbox}[1]{
  \noindent\fbox{%
    \begin{minipage}{0.95\linewidth}
      #1
      \vspace{0.18in}

      \longblank

      \vspace{0.18in}
      \longblank

      \vspace{0.18in}
      \longblank
    \end{minipage}
  }
}

\newcommand{\callout}[3]{
  \noindent\colorbox{#1}{%
    \makebox[\dimexpr0.95\linewidth-2\fboxsep\relax][l]{\textbf{#2}}%
  }

  \noindent\fbox{%
    \begin{minipage}{0.93\linewidth}
      #3
    \end{minipage}
  }
}

\begin{document}

\begin{center}
  {\LARGE \coursenumber{} Week \weekplaceholder{}: Lecture \lectureplaceholder{}}\\
  {\large \lecturetitle{}}
\end{center}

\begin{enumerate}[label=\arabic*.,leftmargin=*,itemsep=.7em]
  \item \textbf{Your name:} \nameblank \hfill \textbf{BU ID:} \idblank
\end{enumerate}

\textbf{Big idea:} A school district wants to know whether new reading software improved student reading levels, but there is concern about a lurking variable (home computer access). We will decide, individually and then with a partner, how the district should sample students to answer this question.

\section*{The Case Study}

A school district purchases language arts software to supplement the teaching of reading to students in grade 2. After 3 months of using the software, the school wants to determine whether use of the program has led to improved reading levels. There is concern that the program may not be as effective for students who are not accustomed to using a computer at home. The district is considering two sampling techniques:

\begin{enumerate}
  \item Randomly sampling students who have a computer at home and students who don't, and assessing their reading skills.
  \item Randomly sampling students regardless of whether they have a computer at home or not, and assessing their reading skills.
\end{enumerate}

\textbf{What is/are the population(s) of interest? Which sampling technique do you think they should use? Why?}

\section*{Step 1: Think (2 mins)}

On your own, consider the question above.

\textbf{My initial thoughts:}
\vfill
\newpage

\section*{Step 2: Pair (3 mins)}

\begin{enumerate}[label=\arabic*.,leftmargin=*,itemsep=.7em, start=2]
  \item \textbf{Partner's name:} \nameblank \hfill \textbf{BU ID:} \idblank
\end{enumerate}

Discuss the case study with your partner. Use these questions to guide your conversation:

\begin{itemize}
  \item What is/are the population(s) of interest in this study?
  \item Which sampling technique do you think the district should use? Why?
  \item Could sampling technique 2 still let home computer access bias the results? How?
\end{itemize}

\textbf{Our thoughts:}
\vfill

\end{document}
